\documentclass[11pt,a4paper,fontset=fandol]{ctexart}

\usepackage[a4paper,top=2.5cm,bottom=2.6cm,left=2.35cm,right=2.35cm,headheight=18pt,headsep=0.55cm,footskip=26pt]{geometry}
\usepackage{amsmath,amssymb,amsthm}
\usepackage{fancyhdr}
\usepackage{lastpage}
\usepackage{multicol}
\usepackage{pdflscape}
\usepackage{enumitem}
\usepackage{titlesec}
\usepackage{xcolor}
\usepackage{hyperref}
\usepackage{bookmark}
\usepackage[most]{tcolorbox}
\usepackage{setspace}
\usepackage{array}

\hypersetup{
  colorlinks=true,
  linkcolor=blue!50!black,
  urlcolor=blue!50!black,
  citecolor=blue!50!black
}

\setstretch{1.18}
\setlength{\parindent}{2em}
\setlength{\parskip}{0.35em}
\setlength{\columnsep}{0.95cm}
\allowdisplaybreaks
\setlist[enumerate]{itemsep=0.3em, topsep=0.3em, leftmargin=2.3em}
\setlist[itemize]{itemsep=0.25em, topsep=0.25em, leftmargin=2em}
\setcounter{tocdepth}{2}

\definecolor{mainblue}{RGB}{36,83,140}
\definecolor{lightblue}{RGB}{241,246,252}
\definecolor{lightgray}{RGB}{246,246,246}
\definecolor{accent}{RGB}{153,76,0}
\definecolor{rulegray}{RGB}{110,118,128}

\titleformat{\section}
  {\Large\bfseries\color{mainblue}}
  {\thesection.}
  {0.45em}
  {}
  [\vspace{0.25em}\color{rulegray}\titlerule]
\titlespacing*{\section}{0pt}{1.3em}{0.9em}
\titleformat{\subsection}{\large\bfseries\color{mainblue}}{\thesubsection}{0.5em}{}
\titlespacing*{\subsection}{0pt}{1.1em}{0.55em}
\titleformat{\subsubsection}{\normalsize\bfseries\color{accent}}{\thesubsubsection}{0.5em}{}

\pagestyle{fancy}
\fancyhf{}
\fancyhead[L]{\small 数论专题课堂讲义}
\fancyhead[C]{\small 同余、一次同余方程组、勾股数与不定方程}
\fancyhead[R]{\small 刘文博}
\fancyfoot[C]{\small 第 \thepage 页 / 共 \pageref{LastPage} 页}
\renewcommand{\headrulewidth}{0.55pt}
\renewcommand{\footrulewidth}{0.35pt}

\newtheoremstyle{formaltheorem}
  {0.8em}{0.8em}{\normalfont}{}{\bfseries\color{mainblue}}{．}{0.6em}{}
\newtheoremstyle{formalexample}
  {0.8em}{0.8em}{\normalfont}{}{\bfseries\color{accent}}{．}{0.6em}{}

\theoremstyle{formaltheorem}
\newtheorem{theorem}{定理}[section]
\newtheorem{proposition}{命题}[section]
\theoremstyle{formalexample}
\newtheorem{example}{例题}[section]
\newtheorem{exercise}{练习}[section]

\newenvironment{solution}{\par\noindent\textbf{解：}}{\hfill$\square$\par}

\newtcolorbox{goalbox}[1]{
  breakable,
  colback=lightblue,
  colframe=mainblue,
  boxrule=0.7pt,
  arc=1mm,
  left=1.2mm,
  right=1.2mm,
  top=1mm,
  bottom=1mm,
  title=\textbf{#1},
  fonttitle=\bfseries
}

\newtcolorbox{notebox}[1]{
  breakable,
  colback=lightgray,
  colframe=black!50,
  boxrule=0.55pt,
  arc=1mm,
  left=1.2mm,
  right=1.2mm,
  top=1mm,
  bottom=1mm,
  title=\textbf{#1},
  fonttitle=\bfseries
}

\newcommand{\classblank}{\vspace{2.2cm}}
\newcommand{\smallblank}{\vspace{1.2cm}}
\newcommand{\ans}[1]{\textbf{答案：}#1}
\newcommand{\partdivider}[2]{%
  \clearpage
  \thispagestyle{empty}
  \vspace*{\fill}
  \begin{center}
  \begin{tcolorbox}[
    enhanced,
    width=0.78\textwidth,
    colback=lightblue,
    colframe=mainblue,
    boxrule=0.9pt,
    arc=0mm,
    left=6mm,
    right=6mm,
    top=8mm,
    bottom=8mm
  ]
  \centering
  {\fontsize{24}{30}\selectfont\bfseries #1\par}
  \vspace{0.6cm}
  {\large #2\par}
  \end{tcolorbox}
  \end{center}
  \vspace*{\fill}
  \clearpage
}

\begin{document}

% =========================
% 封面
% =========================
\thispagestyle{empty}
\begin{titlepage}
\centering
\vspace*{1.2cm}

\begin{tcolorbox}[
  enhanced,
  width=0.92\textwidth,
  colback=white,
  colframe=mainblue,
  boxrule=1pt,
  arc=0mm,
  left=6mm,
  right=6mm,
  top=8mm,
  bottom=8mm
]
\centering

{\fontsize{24}{30}\selectfont\bfseries 数论专题完整课堂讲义\par}
\vspace{0.7cm}
{\fontsize{17}{22}\selectfont 同余、一次同余方程组、中国剩余定理、勾股数与不定方程\par}

\vspace{1.2cm}
\color{rulegray}\rule{0.72\textwidth}{0.7pt}\par
\vspace{1.1cm}

\begin{tcolorbox}[colback=lightblue,colframe=mainblue,width=0.82\textwidth,boxrule=1pt]
\centering
{\Large 适合对象：小学高年级数学竞赛学生}\\[0.4em]
{\large 具备较好的代数基础，接触过方程与因式分解即可}
\end{tcolorbox}

\vspace{1.5cm}

\renewcommand{\arraystretch}{1.4}
\begin{tabular}{>{\bfseries}p{3.5cm}p{8cm}}
讲义作者： & 刘文博 \\
内容结构： & 基础内容与拓展内容结合，适合课堂讲授与课后复习 \\
建议课时： & 4--6 课时，可按学生情况灵活调整 \\
编译方式： & XeLaTeX \\
\end{tabular}

\vspace{1.4cm}
\color{rulegray}\rule{0.72\textwidth}{0.7pt}\par
\vspace{0.9cm}

{\large \bfseries 使用说明\par}
\vspace{0.4cm}
{\normalsize 例题适合课堂讲解与板书示范，练习可先独立完成；文末附录统一给出参考答案，便于课后核对与订正。\par}

{\large \today\par}
\end{tcolorbox}
\end{titlepage}

\pagenumbering{Roman}
\pagestyle{plain}
\tableofcontents
\newpage
\pagestyle{fancy}
\pagenumbering{arabic}

\section{使用说明与学习目标}

\begin{goalbox}{本讲义的使用方式}
本讲义按照由浅入深的顺序编排，分为两个部分：
\begin{itemize}
  \item \textbf{基础篇}：整除、最大公因数、同余、一次不定方程；
  \item \textbf{拓展篇}：一次同余方程组、中国剩余定理、勾股数、进阶不定方程。
\end{itemize}
建议课堂上先完成例题讲解，再安排对应练习；学生可先独立思考，再参考附录答案进行订正。
\end{goalbox}

\begin{goalbox}{本讲目标}
通过本讲的学习，学生将逐步达到以下目标：
\begin{itemize}
  \item 理解整除、最大公因数、欧几里得算法与裴蜀定理；
  \item 学会使用同余处理余数、整除、末位数与不可能性证明；
  \item 会解一次同余方程
  \[
  ax\equiv b\pmod m;
  \]
  \item 会解简单的一次同余方程组，能处理韩信点兵类问题；
  \item 会解二元一次不定方程
  \[
  ax+by=c;
  \]
  \item 会用因式分解法解决平方差型、倒数型、乘积型不定方程；
  \item 了解勾股数的结构与参数表示。
\end{itemize}
\end{goalbox}

\begin{notebox}{授课安排参考}
\begin{itemize}
  \item 第 1 课时：整除、最大公因数、同余基础；
  \item 第 2 课时：同余应用、一次同余方程；
  \item 第 3 课时：一次不定方程、因式分解法；
  \item 第 4 课时：中国剩余定理与韩信点兵；
  \item 第 5 课时：勾股数与进阶不定方程；
  \item 第 6 课时：专题训练与讲评。
\end{itemize}
\end{notebox}

\partdivider{基础篇}{整除、同余与一次不定方程的核心方法}

\section{基础篇：整除与最大公因数}

\subsection{整除}

如果整数 $a,b$ 满足存在整数 $k$ 使得
\[
a=bk,
\]
那么称 $b$ \textbf{整除} $a$，记作
\[
b\mid a.
\]

例如：
\[
3\mid 18,\qquad 5\nmid 18.
\]

\subsection{最大公因数}

两个整数 $a,b$ 的最大公因数记作
\[
(a,b)\quad\text{或}\quad \gcd(a,b).
\]

例如：
\[
(18,24)=6.
\]

\subsection{欧几里得算法}

若
\[
a=bq+r\qquad (0\le r<b),
\]
则
\[
(a,b)=(b,r).
\]
这就是欧几里得算法（辗转相除法）的核心。

\begin{example}
求 $(252,198)$，并把它写成 $252x+198y$ 的形式。
\end{example}

\begin{solution}
先用辗转相除法：
\[
252=198+54,
\]
\[
198=3\times 54+36,
\]
\[
54=36+18,
\]
\[
36=2\times 18.
\]
因此
\[
(252,198)=18.
\]

再反代：
\[
18=54-36,
\]
\[
36=198-3\times 54,
\]
所以
\[
18=54-(198-3\times 54)=4\times 54-198.
\]
又因为
\[
54=252-198,
\]
所以
\[
18=4(252-198)-198=4\times 252-5\times 198.
\]
故
\[
\boxed{18=4\times 252-5\times 198}.
\]
\end{solution}

\subsection{裴蜀定理}

\begin{theorem}[裴蜀定理]
设 $d=(a,b)$，则存在整数 $x,y$，使得
\[
ax+by=d.
\]
\end{theorem}

它立刻推出一个重要结论：

\begin{theorem}
二元一次不定方程
\[
ax+by=c
\]
有整数解，当且仅当
\[
(a,b)\mid c.
\]
\end{theorem}

\begin{notebox}{课堂笔记区}
请写下：为什么“最大公因数能整除右边”会成为一次不定方程有解的必要条件？
\classblank
\end{notebox}

\section{基础篇：同余}

\subsection{同余的定义}

若整数 $a,b,m$ 满足 $m>0$，并且
\[
m\mid (a-b),
\]
那么称 $a$ 与 $b$ \textbf{模 $m$ 同余}，记作
\[
a\equiv b\pmod m.
\]

这句话等价于：$a$ 和 $b$ 除以 $m$ 的余数相同。

例如：
\[
38\equiv 3\pmod 7,
\]
因为
\[
38-3=35,
\]
而 $7\mid 35$。

\subsection{同余的基本性质}

若
\[
a\equiv b\pmod m,\qquad c\equiv d\pmod m,
\]
则有：
\begin{enumerate}
  \item 加法性质：
  \[
  a+c\equiv b+d\pmod m;
  \]
  \item 减法性质：
  \[
  a-c\equiv b-d\pmod m;
  \]
  \item 乘法性质：
  \[
  ac\equiv bd\pmod m;
  \]
  \item 乘方性质：
  \[
  a^n\equiv b^n\pmod m\qquad (n\in \mathbb{N}^+).
  \]
\end{enumerate}

\subsection{约去一个因子要小心}

在同余里，不是任何因子都能随便约掉。只有当
\[
(c,m)=1
\]
时，若
\[
ac\equiv bc\pmod m,
\]
才可推出
\[
a\equiv b\pmod m.
\]

\subsection{常见整除判定的同余解释}

\begin{itemize}
  \item 被 $3$、$9$ 整除：看各位数字和；
  \item 被 $4$ 整除：看末两位；
  \item 被 $8$ 整除：看末三位；
  \item 被 $11$ 整除：看奇数位和与偶数位和的差。
\end{itemize}

例如，因为
\[
10\equiv 1\pmod 9,
\]
所以十进制数
\[
a_n10^n+\cdots +a_110+a_0
\]
模 $9$ 的余数等于
\[
a_n+\cdots +a_1+a_0
\]
模 $9$ 的余数。

\subsection{同余例题}

\begin{example}
求 $2^{100}$ 除以 $7$ 的余数。
\end{example}

\begin{solution}
先找循环：
\[
2^1\equiv 2\pmod 7,
\]
\[
2^2\equiv 4\pmod 7,
\]
\[
2^3\equiv 8\equiv 1\pmod 7.
\]
因此余数循环节为 $3$。因为
\[
100\div 3=33\text{ 余 }1,
\]
所以
\[
2^{100}\equiv 2\pmod 7.
\]
余数是 \boxed{2}。
\end{solution}

\begin{example}
证明：任意整数的平方除以 $4$ 的余数只能是 $0$ 或 $1$。
\end{example}

\begin{solution}
任意整数 $n$ 不是偶数就是奇数。

若 $n=2k$，则
\[
n^2=4k^2\equiv 0\pmod 4.
\]
若 $n=2k+1$，则
\[
n^2=(2k+1)^2=4k^2+4k+1\equiv 1\pmod 4.
\]
所以平方数模 $4$ 只能余 $0$ 或 $1$。
\end{solution}

\begin{example}
判断是否存在整数 $n$，使得
\[
5\mid (n^2+n+1).
\]
\end{example}

\begin{solution}
只需考察 $n$ 模 $5$ 的五种情况：
\[
n\equiv 0,1,2,3,4\pmod 5.
\]
分别代入：
\[
0^2+0+1\equiv 1,
\]
\[
1^2+1+1\equiv 3,
\]
\[
2^2+2+1=7\equiv 2,
\]
\[
3^2+3+1=13\equiv 3,
\]
\[
4^2+4+1=21\equiv 1\pmod 5.
\]
没有一种情况余 $0$，所以不存在这样的整数 $n$。
\end{solution}

\begin{notebox}{课堂笔记区}
请总结：什么时候用“找循环节”，什么时候用“分类讨论模 $m$ 的各种情况”？
\classblank
\end{notebox}

\section{基础篇：一次同余方程}

\subsection{基本结论}

形如
\[
ax\equiv b\pmod m
\]
的方程叫作 \textbf{一次同余方程}。

\begin{theorem}
一次同余方程
\[
ax\equiv b\pmod m
\]
有解，当且仅当
\[
(a,m)\mid b.
\]
\end{theorem}

若设
\[
d=(a,m),
\]
且
\[
d\mid b,
\]
那么模 $m$ 下恰有 $d$ 个两两不同余的解。

\subsection{什么是逆元}

在解一次同余方程时，我们经常会说“求 $a$ 在模 $m$ 下的逆元”。这里的“逆元”虽然名字听起来有些抽象，但含义其实很直观。

\begin{goalbox}{逆元的直观理解}
在普通算式里，如果
\[
a\cdot \frac1a=1,
\]
那么 $\frac1a$ 就相当于 $a$ 的“乘法倒数”。

在模 $m$ 的世界里，如果某个整数 $u$ 满足
\[
au\equiv 1\pmod m,
\]
那么就称 $u$ 是 $a$ 在模 $m$ 下的\textbf{逆元}。
\end{goalbox}

也就是说，逆元就是“模意义下能够把 $a$ 乘回到 $1$ 的那个数”。

例如：
\[
3\times 5=15\equiv 1\pmod 7,
\]
所以 $5$ 是 $3$ 在模 $7$ 下的逆元。

又如：
\[
7\times 7=49\equiv 1\pmod{12},
\]
所以 $7$ 是 $7$ 在模 $12$ 下的逆元。

\begin{theorem}
整数 $a$ 在模 $m$ 下有逆元，当且仅当
\[
(a,m)=1.
\]
\end{theorem}

这条结论非常重要。它告诉我们：只有当 $a$ 与模数 $m$ 互素时，才可以在同余方程里“除以 $a$”。

\subsection{怎样方便地求逆元}

\begin{goalbox}{求逆元的三种常用方法}
\begin{enumerate}
  \item \textbf{直接试乘法}：模数不大时，直接试 $1,2,3,\dots,m-1$；
  \item \textbf{观察凑成 $1$}：利用熟悉的小乘法表快速判断；
  \item \textbf{用裴蜀定理或辗转相除法}：把
  \[
  ax+my=1
  \]
  解出来，其中 $x$ 就是 $a$ 的一个逆元。
\end{enumerate}
\end{goalbox}

\subsubsection{方法一：直接试乘法}

如果模数比较小，最方便的办法就是直接试。

例如要求 $3$ 在模 $7$ 下的逆元，只需试：
\[
3\times 1\equiv 3,\quad
3\times 2\equiv 6,\quad
3\times 3\equiv 2,\quad
3\times 4\equiv 5,\quad
3\times 5\equiv 1\pmod 7.
\]
所以逆元就是
\[
\boxed{5}.
\]

\subsubsection{方法二：利用熟悉恒等变形}

有时可以直接观察。

例如要求 $8$ 在模 $15$ 下的逆元，只需看
\[
8\times 2=16\equiv 1\pmod{15},
\]
所以逆元是
\[
\boxed{2}.
\]

再如要求 $11$ 在模 $12$ 下的逆元，由于
\[
11\equiv -1\pmod{12},
\]
所以
\[
11\times 11\equiv (-1)\times(-1)=1\pmod{12},
\]
故逆元是
\[
\boxed{11}.
\]

\subsubsection{方法三：用裴蜀定理求逆元}

如果 $a$ 与 $m$ 互素，那么根据裴蜀定理，存在整数 $x,y$ 使得
\[
ax+my=1.
\]
把这个式子模 $m$ 来看，就得到
\[
ax\equiv 1\pmod m,
\]
所以 $x$ 就是 $a$ 在模 $m$ 下的一个逆元。

\begin{example}
求 $17$ 在模 $43$ 下的逆元。
\end{example}

\begin{solution}
用辗转相除法：
\[
43=2\times 17+9,
\]
\[
17=1\times 9+8,
\]
\[
9=1\times 8+1.
\]

反代：
\[
1=9-8,
\]
\[
8=17-9,
\]
所以
\[
1=9-(17-9)=2\times 9-17.
\]
又因为
\[
9=43-2\times 17,
\]
所以
\[
1=2(43-2\times 17)-17=2\times 43-5\times 17.
\]
即
\[
(-5)\times 17+2\times 43=1.
\]

因此
\[
17\times(-5)\equiv 1\pmod{43}.
\]
所以 $17$ 在模 $43$ 下的逆元可以取为
\[
\boxed{-5}.
\]
通常把它化成最小正整数，得到
\[
\boxed{38}.
\]
\end{solution}

\begin{notebox}{方法提示}
\begin{itemize}
  \item 逆元不是“普通分数意义上的倒数”，而是“模意义下的乘法倒数”；
  \item 求逆元前一定先检查是否互素；
  \item 小模数优先试乘法，大模数优先用辗转相除法。
\end{itemize}
\end{notebox}

\subsection{标准步骤}

\begin{goalbox}{一次同余方程四步法}
\begin{enumerate}
  \item 求 $d=(a,m)$；
  \item 判断 $d\mid b$ 是否成立；
  \item 若不成立，则无解；
  \item 若成立，则两边除以 $d$，再求逆元。
\end{enumerate}
\end{goalbox}

\begin{example}
解同余方程
\[
7x\equiv 5\pmod{12}.
\]
\end{example}

\begin{solution}
因为
\[
7\times 7=49\equiv 1\pmod{12},
\]
所以 $7$ 在模 $12$ 下的逆元是 $7$。两边同乘 $7$：
\[
x\equiv 35\equiv 11\pmod{12}.
\]
因此
\[
\boxed{x\equiv 11\pmod{12}}.
\]
\end{solution}

\begin{example}
解同余方程
\[
8x\equiv 14\pmod{30}.
\]
\end{example}

\begin{solution}
先求
\[
(8,30)=2.
\]
因为 $2\mid 14$，所以有解。两边除以 $2$，得
\[
4x\equiv 7\pmod{15}.
\]
又因为
\[
4\times 4=16\equiv 1\pmod{15},
\]
所以
\[
x\equiv 4\times 7=28\equiv 13\pmod{15}.
\]
因此模 $30$ 下的全部解为
\[
\boxed{x\equiv 13,\,28\pmod{30}}.
\]
\end{solution}

\begin{example}
判断同余方程
\[
6x\equiv 5\pmod{21}
\]
是否有解。
\end{example}

\begin{solution}
因为
\[
(6,21)=3,
\]
但
\[
3\nmid 5,
\]
所以无解。
\end{solution}

\section{基础篇：二元一次不定方程}

\subsection{基本结论}

\begin{theorem}
方程
\[
ax+by=c
\]
有整数解，当且仅当
\[
(a,b)\mid c.
\]
\end{theorem}

若 $(x_0,y_0)$ 是一组整数解，设
\[
d=(a,b),
\]
则全部整数解为
\[
x=x_0+\frac{b}{d}t,\qquad y=y_0-\frac{a}{d}t,\qquad t\in \mathbb{Z}.
\]

\subsection{例题}

\begin{example}
求方程
\[
12x+18y=84
\]
的所有正整数解。
\end{example}

\begin{solution}
两边同除以 $6$：
\[
2x+3y=14.
\]
容易看出一组解是
\[
x=1,\qquad y=4.
\]
所以全部整数解为
\[
x=1+3t,\qquad y=4-2t,\qquad t\in \mathbb{Z}.
\]
要求正整数解：
\[
1+3t>0,\qquad 4-2t>0.
\]
解得
\[
t=0,1.
\]
故两组正整数解为
\[
\boxed{(x,y)=(1,4),(4,2)}.
\]
\end{solution}

\begin{example}
求方程
\[
5x+7y=59
\]
的所有正整数解。
\end{example}

\begin{solution}
由
\[
5x+7y=59
\]
可得
\[
5x\equiv 59\equiv 3\pmod 7.
\]
因为
\[
5\times 3=15\equiv 1\pmod 7,
\]
所以
\[
x\equiv 3\times 3=9\equiv 2\pmod 7.
\]
设
\[
x=2+7t.
\]
代回原方程：
\[
5(2+7t)+7y=59,
\]
所以
\[
y=7-5t.
\]
要求正整数解，则 $t=0,1$。故正整数解为
\[
\boxed{(x,y)=(2,7),(9,2)}.
\]
\end{solution}

\begin{notebox}{课堂笔记区}
请总结：为什么一次不定方程常常配合同余来做？
\classblank
\end{notebox}

\section{基础篇：因式分解法求不定方程}

\subsection{平方差型}

看到
\[
x^2-y^2=n
\]
时，应立刻想到
\[
x^2-y^2=(x-y)(x+y).
\]

\begin{example}
求正整数解
\[
x^2-y^2=45.
\]
\end{example}

\begin{solution}
化为
\[
(x-y)(x+y)=45.
\]
45 的正因数对有
\[
1\times 45,\quad 3\times 15,\quad 5\times 9.
\]
分别得到
\[
(x,y)=(23,22),(9,6),(7,2).
\]
因此全部正整数解为
\[
\boxed{(23,22),(9,6),(7,2)}.
\]
\end{solution}

\subsection{倒数型}

看到
\[
\frac1x+\frac1y=\frac1n
\]
时，通常通分后配方：
\[
(x-n)(y-n)=n^2.
\]

\begin{example}
求正整数解
\[
\frac1x+\frac1y=\frac16.
\]
\end{example}

\begin{solution}
原方程化为
\[
6x+6y=xy.
\]
移项并配方：
\[
xy-6x-6y+36=36,
\]
即
\[
(x-6)(y-6)=36.
\]
36 的正因数有
\[
1,2,3,4,6,9,12,18,36.
\]
因此全部正整数解为
\[
\boxed{(7,42),(8,24),(9,18),(10,15),(12,12),(15,10),(18,9),(24,8),(42,7)}.
\]
\end{solution}

\subsection{乘积型}

看到
\[
xy=x+y+n
\]
时，常移项并加 $1$：
\[
xy-x-y=n\Rightarrow (x-1)(y-1)=n+1.
\]

\begin{example}
求整数解
\[
x+y+xy=19.
\]
\end{example}

\begin{solution}
两边加 $1$：
\[
xy+x+y+1=20,
\]
即
\[
(x+1)(y+1)=20.
\]
枚举 $20$ 的正负因数对，可得所有整数解：
\[
\boxed{(0,19),(1,9),(3,4),(4,3),(9,1),(19,0),}
\]
\[
\boxed{(-2,-21),(-3,-11),(-5,-6),(-6,-5),(-11,-3),(-21,-2).}
\]
\end{solution}

\partdivider{拓展篇}{一次同余方程组、中国剩余定理、勾股数与进阶不定方程}

\section{拓展篇：一次同余方程组}

\subsection{最朴素的方法：逐步代入}

\begin{example}
解方程组
\[
\begin{cases}
x\equiv 2\pmod 3,\\
x\equiv 3\pmod 5.
\end{cases}
\]
\end{example}

\begin{solution}
由第一式设
\[
x=2+3t.
\]
代入第二式：
\[
2+3t\equiv 3\pmod 5,
\]
即
\[
3t\equiv 1\pmod 5.
\]
因为 $3$ 在模 $5$ 下的逆元是 $2$，所以
\[
t\equiv 2\pmod 5.
\]
设
\[
t=2+5k,
\]
则
\[
x=2+3(2+5k)=8+15k.
\]
所以
\[
\boxed{x\equiv 8\pmod{15}}.
\]
\end{solution}

\subsection{模数不互素时的判定}

\begin{theorem}
同余方程组
\[
\begin{cases}
x\equiv a\pmod m,\\
x\equiv b\pmod n
\end{cases}
\]
有解，当且仅当
\[
a\equiv b\pmod{(m,n)}.
\]
若有解，则解在模 $[m,n]$ 的意义下唯一。
\end{theorem}

\begin{example}
解方程组
\[
\begin{cases}
x\equiv 1\pmod 4,\\
x\equiv 3\pmod 6.
\end{cases}
\]
\end{example}

\begin{solution}
因为
\[
(4,6)=2,
\]
且
\[
1\equiv 3\pmod 2,
\]
所以有解。

由第一式设
\[
x=1+4t.
\]
代入第二式：
\[
1+4t\equiv 3\pmod 6,
\]
即
\[
4t\equiv 2\pmod 6.
\]
两边除以 $2$：
\[
2t\equiv 1\pmod 3.
\]
所以
\[
t\equiv 2\pmod 3.
\]
设
\[
t=2+3k,
\]
则
\[
x=9+12k.
\]
因此
\[
\boxed{x\equiv 9\pmod{12}}.
\]
\end{solution}

\section{拓展篇：中国剩余定理}

\begin{theorem}[中国剩余定理]
如果
\[
m_1,m_2,\dots,m_k
\]
两两互素，那么同余方程组
\[
\begin{cases}
x\equiv a_1\pmod{m_1},\\
x\equiv a_2\pmod{m_2},\\
\cdots\\
x\equiv a_k\pmod{m_k}
\end{cases}
\]
一定有解，而且解在模
\[
M=m_1m_2\cdots m_k
\]
的意义下唯一。
\end{theorem}

\begin{goalbox}{一句话理解}
当模数两两互素时，多个余数条件不会打架，而且一定能拼成一个唯一的总答案。
\end{goalbox}

\subsection{韩信点兵}

\begin{example}[韩信点兵]
某数被 $3$ 除余 $2$，被 $5$ 除余 $3$，被 $7$ 除余 $2$。求这个数最小是多少？并写出全部解。
\end{example}

\begin{solution}
要求解
\[
\begin{cases}
x\equiv 2\pmod 3,\\
x\equiv 3\pmod 5,\\
x\equiv 2\pmod 7.
\end{cases}
\]

由第一式设
\[
x=2+3t.
\]
代入第二式：
\[
2+3t\equiv 3\pmod 5
\Rightarrow 3t\equiv 1\pmod 5.
\]
所以
\[
t\equiv 2\pmod 5.
\]
设
\[
t=2+5s,
\]
则
\[
x=8+15s.
\]
再代入第三式：
\[
8+15s\equiv 2\pmod 7
\Rightarrow 1+s\equiv 2\pmod 7.
\]
所以
\[
s\equiv 1\pmod 7.
\]
设
\[
s=1+7k,
\]
则
\[
x=23+105k.
\]
因此最小正整数解为
\[
\boxed{23},
\]
全部解为
\[
\boxed{x\equiv 23\pmod{105}}.
\]
\end{solution}

\subsection{构造法简介}

设
\[
M=m_1m_2\cdots m_k,\qquad M_i=\frac{M}{m_i}.
\]
因为
\[
(M_i,m_i)=1,
\]
所以存在整数 $t_i$ 使得
\[
M_it_i\equiv 1\pmod{m_i}.
\]
那么
\[
x\equiv a_1M_1t_1+a_2M_2t_2+\cdots+a_kM_kt_k\pmod M
\]
就是一个解。

\begin{example}
解方程组
\[
\begin{cases}
x\equiv 1\pmod 2,\\
x\equiv 2\pmod 3,\\
x\equiv 4\pmod 5.
\end{cases}
\]
\end{example}

\begin{solution}
总模数
\[
M=2\times 3\times 5=30.
\]
取
\[
M_1=15,\qquad M_2=10,\qquad M_3=6.
\]
又因为
\[
15\equiv 1\pmod 2,\qquad 10\equiv 1\pmod 3,\qquad 6\equiv 1\pmod 5,
\]
所以三个逆元都可取为 $1$。
因此
\[
x\equiv 1\cdot 15+2\cdot 10+4\cdot 6=59\equiv 29\pmod{30}.
\]
故
\[
\boxed{x\equiv 29\pmod{30}}.
\]
\end{solution}

\begin{notebox}{课堂笔记区}
请自己写一个“韩信点兵题”的解题模板：
\classblank
\end{notebox}

\section{拓展篇：勾股数}

\subsection{基本概念}

满足
\[
a^2+b^2=c^2
\]
的正整数三元组 $(a,b,c)$ 叫作 \textbf{勾股数}。

例如：
\[
(3,4,5),\quad (5,12,13),\quad (8,15,17).
\]

若
\[
(a,b,c)=1,
\]
则称为 \textbf{原始勾股数}。

\begin{proposition}
若 $(a,b,c)$ 是原始勾股数，则：
\begin{enumerate}
  \item $a,b$ 中恰有一个是偶数；
  \item $c$ 一定是奇数；
  \item $a,b$ 互素。
\end{enumerate}
\end{proposition}

\begin{proof}
因为 $(a,b,c)$ 是原始勾股数，所以
\[
a^2+b^2=c^2,
\qquad (a,b,c)=1.
\]

先证明 $a,b$ 不可能同时为偶数。若 $a,b$ 都是偶数，那么 $a^2,b^2$ 都是偶数，
从而 $c^2=a^2+b^2$ 也是偶数，于是 $c$ 也是偶数。这样一来，$a,b,c$ 都能被 $2$
整除，这与
\[
(a,b,c)=1
\]
矛盾。

再证明 $a,b$ 不可能同时为奇数。若 $a,b$ 都是奇数，则奇数平方模 $4$ 都余 $1$，
所以
\[
a^2\equiv 1\pmod 4,\qquad b^2\equiv 1\pmod 4.
\]
因此
\[
c^2=a^2+b^2\equiv 1+1\equiv 2\pmod 4.
\]
但任意整数的平方模 $4$ 只能余 $0$ 或 $1$，不可能余 $2$，矛盾。

由于 $a,b$ 既不可能同时为偶数，也不可能同时为奇数，所以 $a,b$ 中恰有一个是偶数，
另一个是奇数，这就证明了第 1 条。

由第 1 条可知，$a^2,b^2$ 中恰有一个是偶数，另一个是奇数，所以
\[
c^2=a^2+b^2
\]
是奇数，从而 $c$ 也是奇数，这就证明了第 2 条。

最后证明 $a,b$ 互素。若 $a,b$ 不互素，那么存在整数 $d>1$，使得
\[
d\mid a,\qquad d\mid b.
\]
于是
\[
d\mid a^2,\qquad d\mid b^2,
\]
从而
\[
d\mid (a^2+b^2)=c^2.
\]
这说明 $d$ 是 $c^2$ 的因数。由于素因数若整除一个平方数，就必整除这个数本身，
所以 $d$ 的每个素因数都整除 $c$，从而 $a,b,c$ 有公共素因数，这与
\[
(a,b,c)=1
\]
矛盾。因此 $a,b$ 互素，第 3 条也得证。
\end{proof}

\subsection{参数公式}

\begin{theorem}
所有原始勾股数都可以表示为
\[
a=m^2-n^2,\qquad b=2mn,\qquad c=m^2+n^2,
\]
其中
\[
m>n\ge 1,\qquad (m,n)=1,
\]
并且 $m,n$ 一奇一偶。
\end{theorem}

\begin{proof}
设 $(a,b,c)$ 是一个原始勾股数，即
\[
a^2+b^2=c^2,
\qquad (a,b,c)=1.
\]
由前面的命题可知，$a,b$ 中恰有一个是偶数，$c$ 是奇数。不妨设
\[
b \text{ 为偶数},\qquad a \text{ 为奇数}.
\]

由
\[
a^2+b^2=c^2
\]
可得
\[
b^2=c^2-a^2=(c+a)(c-a).
\]
因为 $a,c$ 都是奇数，所以
\[
c+a,\qquad c-a
\]
都是偶数。于是可设
\[
c+a=2u,\qquad c-a=2v,
\]
其中 $u,v$ 是正整数，且 $u>v$。将它们代回上式，得到
\[
b^2=(2u)(2v)=4uv,
\]
所以
\[
\left(\frac{b}{2}\right)^2=uv.
\]

下面证明 $u,v$ 互素。若整数 $d$ 同时整除 $u,v$，则 $d$ 整除它们的和与差：
\[
d\mid (u+v)=c,\qquad d\mid (u-v)=a.
\]
于是 $d$ 同时整除 $a,c$。再由
\[
b^2=c^2-a^2
\]
可知 $d$ 也整除 $b^2$，从而整除 $b$。这就说明 $a,b,c$ 有公共因数，与
\[
(a,b,c)=1
\]
矛盾。因此
\[
(u,v)=1.
\]

现在，$u$ 与 $v$ 互素，并且它们的乘积 $uv$ 是一个完全平方数。于是 $u,v$
本身都必须是完全平方数。设
\[
u=m^2,\qquad v=n^2,
\]
其中
\[
m>n\ge 1.
\]
那么
\[
\left(\frac{b}{2}\right)^2=uv=m^2n^2,
\]
所以
\[
\frac{b}{2}=mn,
\]
即
\[
b=2mn.
\]

再由
\[
c+a=2m^2,\qquad c-a=2n^2
\]
两式相加、相减，得到
\[
c=m^2+n^2,\qquad a=m^2-n^2.
\]
于是
\[
a=m^2-n^2,\qquad b=2mn,\qquad c=m^2+n^2.
\]

下面证明 $(m,n)=1$ 且 $m,n$ 一奇一偶。

若 $m,n$ 有公因数 $d>1$，则 $d^2$ 同时整除 $m^2$ 与 $n^2$，从而整除
\[
a=m^2-n^2,\qquad c=m^2+n^2.
\]
又因为
\[
b=2mn,
\]
所以 $d$ 也整除 $b$。这与 $(a,b,c)=1$ 矛盾。因此
\[
(m,n)=1.
\]

再看奇偶性。若 $m,n$ 同为偶数，则它们不互素；若 $m,n$ 同为奇数，则
\[
a=m^2-n^2
\]
为偶数，这与前面约定的 $a$ 为奇数矛盾。因此 $m,n$ 必为一奇一偶。

这就证明了：任意原始勾股数都可以表示成
\[
a=m^2-n^2,\qquad b=2mn,\qquad c=m^2+n^2,
\]
其中
\[
m>n\ge 1,\qquad (m,n)=1,
\]
并且 $m,n$ 一奇一偶。

反过来，若 $m,n$ 满足这些条件，令
\[
a=m^2-n^2,\qquad b=2mn,\qquad c=m^2+n^2,
\]
则直接计算可得
\[
a^2+b^2=(m^2-n^2)^2+(2mn)^2=(m^2+n^2)^2=c^2.
\]
因此 $(a,b,c)$ 是勾股数。

再证明它是原始的。设 $d$ 同时整除 $a,b$，则
\[
d\mid (m^2-n^2),\qquad d\mid 2mn.
\]
由于 $a=m^2-n^2$ 为奇数，所以 $d$ 也是奇数。于是若素数 $p$ 整除 $d$，则
$p$ 整除 $mn$，不妨设 $p\mid m$。又因为
\[
p\mid (m^2-n^2)
\]
且 $p\mid m^2$，所以 $p\mid n^2$，从而 $p\mid n$。这与
\[
(m,n)=1
\]
矛盾。因此 $(a,b)=1$。

再设某个整数 $d$ 同时整除 $a,c$。由于 $a,c$ 都是奇数，所以 $d$ 也是奇数。
若素数 $p$ 整除 $d$，则
\[
p\mid (a+c)=2m^2,\qquad p\mid (c-a)=2n^2.
\]
因为 $p$ 是奇素数，所以
\[
p\mid m^2,\qquad p\mid n^2,
\]
从而
\[
p\mid m,\qquad p\mid n,
\]
这与 $(m,n)=1$ 矛盾。因此 $(a,c)=1$。

同理，若某个整数 $d$ 同时整除 $b,c$，则任何整除 $d$ 的素数 $p$ 都满足
\[
p\mid 2mn,\qquad p\mid m^2+n^2.
\]
因为 $c=m^2+n^2$ 为奇数，所以 $p\ne 2$。不妨设 $p\mid m$，则由
\[
p\mid (m^2+n^2)
\]
可得 $p\mid n^2$，从而 $p\mid n$，仍与 $(m,n)=1$ 矛盾。因此 $(b,c)=1$。

于是
\[
(a,b)=(a,c)=(b,c)=1,
\]
从而
\[
(a,b,c)=1.
\]
所以由这样的 $m,n$ 构造出的勾股数一定是原始勾股数。
\end{proof}

\begin{example}
写出斜边小于 $30$ 的全部原始勾股数。
\end{example}

\begin{solution}
枚举
\[
c=m^2+n^2<30.
\]
满足条件的 $(m,n)$ 为：
\begin{enumerate}
  \item $m=2,n=1$，得 $(3,4,5)$；
  \item $m=3,n=2$，得 $(5,12,13)$；
  \item $m=4,n=1$，得 $(8,15,17)$；
  \item $m=4,n=3$，得 $(7,24,25)$；
  \item $m=5,n=2$，得 $(20,21,29)$。
\end{enumerate}
因此全部原始勾股数为
\[
\boxed{(3,4,5),(5,12,13),(8,15,17),(7,24,25),(20,21,29)}.
\]
\end{solution}

\begin{example}
求周长为 $84$ 的全部直角三角形整数边长。
\end{example}

\begin{solution}
设三边为
\[
ka,\ kb,\ kc,
\]
其中 $(a,b,c)$ 是原始勾股数，则
\[
k(a+b+c)=84.
\]
检查常见原始勾股数的周长：
\begin{itemize}
  \item $(3,4,5)$ 的周长为 $12$，可取 $k=7$，得 $(21,28,35)$；
  \item $(12,35,37)$ 的周长正好为 $84$。
\end{itemize}
因此全部整数边长为
\[
\boxed{(21,28,35),(12,35,37)}.
\]
\end{solution}

\begin{notebox}{课堂笔记区}
请写下你最熟悉的 5 组常见勾股数：
\classblank
\end{notebox}

\section{拓展篇：进阶不定方程}

\subsection{平方差型进阶}

\begin{theorem}
方程
\[
x^2-y^2=n
\]
有整数解，当且仅当
\[
n\not\equiv 2\pmod 4.
\]
\end{theorem}

\begin{example}
求正整数解
\[
x^2-y^2=105.
\]
\end{example}

\begin{solution}
化为
\[
(x-y)(x+y)=105.
\]
枚举正因数对：
\[
1\times 105,\quad 3\times 35,\quad 5\times 21,\quad 7\times 15.
\]
分别得到
\[
(x,y)=(53,52),(19,16),(13,8),(11,4).
\]
故全部正整数解为
\[
\boxed{(53,52),(19,16),(13,8),(11,4)}.
\]
\end{solution}

\begin{example}
判断方程
\[
x^2-y^2=2026
\]
是否有整数解，并说明理由。
\end{example}

\begin{solution}
因为
\[
2026\equiv 2\pmod 4,
\]
而平方差不可能模 $4$ 余 $2$，所以原方程无整数解。
\end{solution}

\subsection{倒数型进阶}

\begin{example}
求正整数解
\[
\frac1x+\frac1y=\frac1{12}.
\]
\end{example}

\begin{solution}
化为
\[
12x+12y=xy.
\]
移项并配方：
\[
xy-12x-12y+144=144,
\]
所以
\[
(x-12)(y-12)=144.
\]
若设 $d\mid 144$ 且 $d>0$，则
\[
x=12+d,\qquad y=12+\frac{144}{d}.
\]
故全部正整数解可统一写成
\[
\boxed{x=12+d,\quad y=12+\frac{144}{d}\qquad (d\mid 144,\ d>0).}
\]
\end{solution}

\subsection{乘积型进阶}

\begin{example}
求正整数解
\[
xy=x+y+56.
\]
\end{example}

\begin{solution}
移项并加 $1$：
\[
xy-x-y=56\Rightarrow (x-1)(y-1)=57.
\]
因为
\[
57=1\times 57=3\times 19,
\]
所以全部正整数解为
\[
\boxed{(2,58),(4,20),(20,4),(58,2)}.
\]
\end{solution}

\subsection{线性不定方程与同余结合}

\begin{example}
求正整数解
\[
7x+11y=100.
\]
\end{example}

\begin{solution}
由
\[
7x+11y=100
\]
可得
\[
7x\equiv 100\equiv 1\pmod{11}.
\]
因为
\[
7\times 8=56\equiv 1\pmod{11},
\]
所以
\[
x\equiv 8\pmod{11}.
\]
设
\[
x=8+11t,
\]
代回得
\[
y=4-7t.
\]
为使 $y>0$，只能取 $t=0$，故唯一正整数解为
\[
\boxed{(x,y)=(8,4)}.
\]
\end{solution}

\begin{example}
证明方程
\[
x^2+2=y^2
\]
没有整数解。
\end{example}

\begin{solution}
移项得
\[
y^2-x^2=2,
\]
即
\[
(y-x)(y+x)=2.
\]
而 $y-x$ 与 $y+x$ 同奇偶：
\begin{itemize}
  \item 若都为奇数，则乘积为奇数；
  \item 若都为偶数，则乘积至少能被 $4$ 整除。
\end{itemize}
都不可能等于 $2$，故原方程无整数解。
\end{solution}

\section{方法总表}

\begin{goalbox}{看见什么，就想到什么}
\begin{itemize}
  \item 看见大幂次、末位数、余数问题：想到 \textbf{同余}；
  \item 看见一次同余方程：先看 \textbf{$(a,m)\mid b$}；
  \item 看见多个余数条件：想到 \textbf{逐步代入法 / 中国剩余定理}；
  \item 看见一次不定方程：先看 \textbf{$(a,b)\mid c$}；
  \item 看见 $x^2-y^2=n$：想到 \textbf{平方差分解}；
  \item 看见 $\frac1x+\frac1y=\frac1n$：想到 \textbf{$(x-n)(y-n)=n^2$}；
  \item 看见 $xy=x+y+n$：想到 \textbf{$(x-1)(y-1)=n+1$}；
  \item 看见直角三角形整数边：想到 \textbf{勾股数参数公式}。
\end{itemize}
\end{goalbox}

\section{课堂练习}

\begin{goalbox}{练习使用说明}
A 组以基础巩固为主，B 组加入适度拓展内容。建议先独立完成，再翻到附录核对。
\end{goalbox}

\subsection{A组：基础巩固}

\begin{multicols}{2}

\begin{exercise}
求下列各式的余数：
\begin{enumerate}
  \item $37^{50}$ 除以 $8$ 的余数；
  \item $2^{60}+3^{60}$ 除以 $5$ 的余数；
  \item $10^{100}+1$ 除以 $11$ 的余数。
\end{enumerate}
\end{exercise}
\smallblank

\begin{exercise}
解下列同余方程：
\begin{enumerate}
  \item $3x\equiv 4\pmod 7$；
  \item $6x\equiv 9\pmod{15}$。
\end{enumerate}
\end{exercise}
\smallblank

\begin{exercise}
证明：任意整数的平方除以 $8$ 的余数只能是 $0,1,4$。
\end{exercise}
\smallblank

\begin{exercise}
求所有小于 $50$ 的正整数 $n$，满足
\[
n\equiv 2\pmod 3,\qquad n\equiv 1\pmod 4.
\]
\end{exercise}
\smallblank

\begin{exercise}
求所有整数 $n$，使得
\[
7\mid (n^2+3n+4).
\]
\end{exercise}
\smallblank

\begin{exercise}
求最小正整数 $n$，使得
\[
n\equiv 3\pmod 5,\qquad n\equiv 4\pmod 7.
\]
并写出全部整数解。
\end{exercise}
\smallblank

\begin{exercise}
求方程
\[
14x+21y=35
\]
的所有整数解。
\end{exercise}
\smallblank

\begin{exercise}
求方程
\[
3x+4y=100
\]
的所有正整数解。
\end{exercise}
\smallblank

\begin{exercise}
判断方程
\[
15x+25y=7
\]
是否有整数解，并说明理由。
\end{exercise}
\smallblank

\begin{exercise}
求正整数解
\[
x^2-y^2=36.
\]
\end{exercise}
\smallblank

\begin{exercise}
求正整数解
\[
\frac1x+\frac1y=\frac14.
\]
\end{exercise}
\smallblank

\begin{exercise}
求整数解
\[
x+y+xy=19.
\]
\end{exercise}
\smallblank

\end{multicols}

\subsection{B组：拓展训练}

\begin{multicols}{2}

\begin{exercise}
解下列同余方程：
\begin{enumerate}
  \item $10x\equiv 15\pmod{35}$；
  \item $12x\equiv 7\pmod{25}$；
  \item $14x\equiv 5\pmod{21}$。
\end{enumerate}
\end{exercise}
\smallblank

\begin{exercise}
解方程组
\[
\begin{cases}
x\equiv 1\pmod 2,\\
x\equiv 2\pmod 3,\\
x\equiv 4\pmod 5.
\end{cases}
\]
\end{exercise}
\smallblank

\begin{exercise}
解方程组
\[
\begin{cases}
x\equiv 3\pmod 7,\\
x\equiv 4\pmod 9,\\
x\equiv 5\pmod{11}.
\end{cases}
\]
\end{exercise}
\smallblank

\begin{exercise}
解方程组
\[
\begin{cases}
x\equiv 2\pmod 4,\\
x\equiv 8\pmod 6.
\end{cases}
\]
\end{exercise}
\smallblank

\begin{exercise}
判断下列方程组是否有解：
\[
\begin{cases}
x\equiv 1\pmod 6,\\
x\equiv 4\pmod 9.
\end{cases}
\]
\end{exercise}
\smallblank

\begin{exercise}
某数被 $4$ 除余 $3$，被 $5$ 除余 $1$，被 $7$ 除余 $6$。求最小正整数解。
\end{exercise}
\smallblank

\begin{exercise}
写出斜边小于 $40$ 的全部原始勾股数。
\end{exercise}
\smallblank

\begin{exercise}
求周长为 $60$ 的全部直角三角形整数边长。
\end{exercise}
\smallblank

\begin{exercise}
求正整数解
\[
x^2-y^2=96.
\]
\end{exercise}
\smallblank

\begin{exercise}
判断方程
\[
x^2-y^2=2026
\]
是否有整数解，并说明理由。
\end{exercise}
\smallblank

\begin{exercise}
求正整数解
\[
xy=x+y+35.
\]
\end{exercise}
\smallblank

\begin{exercise}
求正整数解
\[
\frac1x+\frac1y=\frac18.
\]
\end{exercise}
\smallblank

\begin{exercise}
求正整数解
\[
9x+14y=121.
\]
\end{exercise}
\smallblank

\begin{exercise}
证明：不存在整数 $x,y$，使得
\[
x^2+y^2=3.
\]
\end{exercise}
\smallblank

\end{multicols}

\begin{notebox}{随堂总结区}
请写出今天你最重要的 3 个收获：
\classblank
\end{notebox}

\newpage
\appendix
\section{附录：练习参考答案}

\subsection{A组参考答案}

\begin{exercise}[A1]
求下列各式的余数：
\begin{enumerate}
  \item $37^{50}$ 除以 $8$ 的余数；
  \item $2^{60}+3^{60}$ 除以 $5$ 的余数；
  \item $10^{100}+1$ 除以 $11$ 的余数。
\end{enumerate}
\end{exercise}

\begin{solution}
\begin{enumerate}
  \item 因为
  \[
  37\equiv 5\pmod 8,\qquad 5^2=25\equiv 1\pmod 8,
  \]
  所以
  \[
  37^{50}\equiv 5^{50}=(5^2)^{25}\equiv 1\pmod 8.
  \]
  余数为 \boxed{1}。

  \item 因为
  \[
  2^4\equiv 1\pmod 5,\qquad 3^4\equiv 1\pmod 5,
  \]
  又 $60$ 是 $4$ 的倍数，所以
  \[
  2^{60}\equiv 1,\qquad 3^{60}\equiv 1\pmod 5.
  \]
  因此
  \[
  2^{60}+3^{60}\equiv 2\pmod 5.
  \]
  余数为 \boxed{2}。

  \item 因为
  \[
  10\equiv -1\pmod{11},
  \]
  所以
  \[
  10^{100}+1\equiv (-1)^{100}+1=2\pmod{11}.
  \]
  余数为 \boxed{2}。
\end{enumerate}
\end{solution}

\begin{exercise}[A2]
解下列同余方程：
\begin{enumerate}
  \item $3x\equiv 4\pmod 7$；
  \item $6x\equiv 9\pmod{15}$。
\end{enumerate}
\end{exercise}

\begin{solution}
\begin{enumerate}
  \item 因为
  \[
  3\times 5=15\equiv 1\pmod 7,
  \]
  所以
  \[
  x\equiv 5\times 4=20\equiv 6\pmod 7.
  \]
  即
  \[
  \boxed{x\equiv 6\pmod 7}.
  \]

  \item 因为
  \[
  (6,15)=3,\qquad 3\mid 9,
  \]
  所以有解。两边除以 $3$ 得
  \[
  2x\equiv 3\pmod 5.
  \]
  又因 $2^{-1}\equiv 3\pmod 5$，故
  \[
  x\equiv 3\times 3=9\equiv 4\pmod 5.
  \]
  因此全部解为
  \[
  \boxed{x=4+5k\quad (k\in \mathbb{Z})}.
  \]
\end{enumerate}
\end{solution}

\begin{exercise}[A3]
证明：任意整数的平方除以 $8$ 的余数只能是 $0,1,4$。
\end{exercise}

\begin{solution}
分奇偶讨论：

若 $n$ 为偶数，设 $n=2k$，则
\[
n^2=4k^2.
\]
若 $k$ 为偶数，则 $n^2\equiv 0\pmod 8$；若 $k$ 为奇数，则 $n^2\equiv 4\pmod 8$。

若 $n$ 为奇数，设 $n=2k+1$，则
\[
n^2=(2k+1)^2=4k(k+1)+1.
\]
由于 $k(k+1)$ 为偶数，所以
\[
n^2\equiv 1\pmod 8.
\]
综上，平方数模 $8$ 只能余
\[
\boxed{0,1,4}.
\]
\end{solution}

\begin{exercise}[A4]
求所有小于 $50$ 的正整数 $n$，满足
\[
n\equiv 2\pmod 3,\qquad n\equiv 1\pmod 4.
\]
\end{exercise}

\begin{solution}
列出小于 $50$ 且模 $3$ 余 $2$ 的数：
\[
2,5,8,11,14,17,20,23,26,29,32,35,38,41,44,47.
\]
其中模 $4$ 余 $1$ 的是
\[
\boxed{5,17,29,41}.
\]
\end{solution}

\begin{exercise}[A5]
求所有整数 $n$，使得
\[
7\mid (n^2+3n+4).
\]
\end{exercise}

\begin{solution}
注意到
\[
4(n^2+3n+4)=(2n+3)^2+7.
\]
若 $7\mid (n^2+3n+4)$，则
\[
(2n+3)^2\equiv 0\pmod 7,
\]
所以
\[
2n+3\equiv 0\pmod 7.
\]
即
\[
2n\equiv 4\pmod 7.
\]
因为 $2^{-1}\equiv 4\pmod 7$，故
\[
n\equiv 16\equiv 2\pmod 7.
\]
因此全部解为
\[
\boxed{n\equiv 2\pmod 7}.
\]
\end{solution}

\begin{exercise}[A6]
求最小正整数 $n$，使得
\[
n\equiv 3\pmod 5,\qquad n\equiv 4\pmod 7.
\]
并写出全部整数解。
\end{exercise}

\begin{solution}
列出模 $5$ 余 $3$ 的数：
\[
3,8,13,18,23,\dots
\]
其中模 $7$ 余 $4$ 的最小正整数是
\[
\boxed{18}.
\]
因为 $5,7$ 互素，所以全部解为
\[
\boxed{n=18+35k\quad (k\in \mathbb{Z})}.
\]
\end{solution}

\begin{exercise}[A7]
求方程
\[
14x+21y=35
\]
的所有整数解。
\end{exercise}

\begin{solution}
两边同除以 $7$，得
\[
2x+3y=5.
\]
一组解为
\[
x=1,\qquad y=1.
\]
故全部整数解为
\[
\boxed{x=1+3t,\qquad y=1-2t\qquad (t\in \mathbb{Z})}.
\]
\end{solution}

\begin{exercise}[A8]
求方程
\[
3x+4y=100
\]
的所有正整数解。
\end{exercise}

\begin{solution}
由
\[
3x\equiv 0\pmod 4
\]
得
\[
x\equiv 0\pmod 4.
\]
设
\[
x=4t.
\]
代回得
\[
y=25-3t.
\]
要求正整数解，所以
\[
t=1,2,3,4,5,6,7,8.
\]
故全部正整数解为
\[
\boxed{(4,22),(8,19),(12,16),(16,13),(20,10),(24,7),(28,4),(32,1)}.
\]
\end{solution}

\begin{exercise}[A9]
判断方程
\[
15x+25y=7
\]
是否有整数解，并说明理由。
\end{exercise}

\begin{solution}
因为
\[
(15,25)=5,
\]
但
\[
5\nmid 7,
\]
所以无整数解。
\end{solution}

\begin{exercise}[A10]
求正整数解
\[
x^2-y^2=36.
\]
\end{exercise}

\begin{solution}
化为
\[
(x-y)(x+y)=36.
\]
在正因数对中，只有
\[
2\times 18
\]
能得到正整数解：
\[
x=10,\qquad y=8.
\]
所以唯一正整数解为
\[
\boxed{(10,8)}.
\]
\end{solution}

\begin{exercise}[A11]
求正整数解
\[
\frac1x+\frac1y=\frac14.
\]
\end{exercise}

\begin{solution}
化为
\[
(x-4)(y-4)=16.
\]
枚举 $16$ 的正因数，可得
\[
\boxed{(5,20),(6,12),(8,8),(12,6),(20,5)}.
\]
\end{solution}

\begin{exercise}[A12]
求整数解
\[
x+y+xy=19.
\]
\end{exercise}

\begin{solution}
化为
\[
(x+1)(y+1)=20.
\]
枚举正负因数对，得全部整数解：
\[
\boxed{(0,19),(1,9),(3,4),(4,3),(9,1),(19,0),}
\]
\[
\boxed{(-2,-21),(-3,-11),(-5,-6),(-6,-5),(-11,-3),(-21,-2).}
\]
\end{solution}

\subsection{B组参考答案}

\begin{exercise}[B1]
解下列同余方程：
\begin{enumerate}
  \item $10x\equiv 15\pmod{35}$；
  \item $12x\equiv 7\pmod{25}$；
  \item $14x\equiv 5\pmod{21}$。
\end{enumerate}
\end{exercise}

\begin{solution}
\begin{enumerate}
  \item 因为
  \[
  (10,35)=5,\qquad 5\mid 15,
  \]
  两边除以 $5$ 得
  \[
  2x\equiv 3\pmod 7.
  \]
  所以
  \[
  x\equiv 5\pmod 7.
  \]
  因此模 $35$ 下的全部解为
  \[
  \boxed{x\equiv 5,12,19,26,33\pmod{35}}.
  \]

  \item 因为
  \[
  (12,25)=1,
  \]
  且
  \[
  12\times 23=276\equiv 1\pmod{25},
  \]
  所以
  \[
  x\equiv 23\times 7=161\equiv 11\pmod{25}.
  \]
  即
  \[
  \boxed{x\equiv 11\pmod{25}}.
  \]

  \item 因为
  \[
  (14,21)=7,
  \]
  但
  \[
  7\nmid 5,
  \]
  所以无解。
\end{enumerate}
\end{solution}

\begin{exercise}[B2]
解方程组
\[
\begin{cases}
x\equiv 1\pmod 2,\\
x\equiv 2\pmod 3,\\
x\equiv 4\pmod 5.
\end{cases}
\]
\end{exercise}

\begin{solution}
由中国剩余定理可得
\[
\boxed{x\equiv 29\pmod{30}}.
\]
\end{solution}

\begin{exercise}[B3]
解方程组
\[
\begin{cases}
x\equiv 3\pmod 7,\\
x\equiv 4\pmod 9,\\
x\equiv 5\pmod{11}.
\end{cases}
\]
\end{exercise}

\begin{solution}
设
\[
x=3+7t.
\]
代入第二式，得
\[
7t\equiv 1\pmod 9\Rightarrow t\equiv 4\pmod 9.
\]
设
\[
t=4+9s,
\]
则
\[
x=31+63s.
\]
代入第三式：
\[
31+63s\equiv 5\pmod{11}
\Rightarrow 8s\equiv 7\pmod{11}.
\]
因为 $8^{-1}\equiv 7\pmod{11}$，所以
\[
s\equiv 5\pmod{11}.
\]
设
\[
s=5+11k,
\]
则
\[
x=346+693k.
\]
故
\[
\boxed{x\equiv 346\pmod{693}}.
\]
\end{solution}

\begin{exercise}[B4]
解方程组
\[
\begin{cases}
x\equiv 2\pmod 4,\\
x\equiv 8\pmod 6.
\end{cases}
\]
\end{exercise}

\begin{solution}
第二式等价于
\[
x\equiv 2\pmod 6.
\]
所以方程组变成
\[
\begin{cases}
x\equiv 2\pmod 4,\\
x\equiv 2\pmod 6.
\end{cases}
\]
故
\[
\boxed{x\equiv 2\pmod{12}}.
\]
\end{solution}

\begin{exercise}[B5]
判断下列方程组是否有解：
\[
\begin{cases}
x\equiv 1\pmod 6,\\
x\equiv 4\pmod 9.
\end{cases}
\]
\end{exercise}

\begin{solution}
因为
\[
(6,9)=3,
\]
且
\[
1\equiv 4\pmod 3,
\]
所以有解。
设
\[
x=1+6t,
\]
代入第二式，得
\[
6t\equiv 3\pmod 9\Rightarrow 2t\equiv 1\pmod 3.
\]
故
\[
t\equiv 2\pmod 3.
\]
所以
\[
\boxed{x\equiv 13\pmod{18}}.
\]
\end{solution}

\begin{exercise}[B6]
某数被 $4$ 除余 $3$，被 $5$ 除余 $1$，被 $7$ 除余 $6$。求最小正整数解。
\end{exercise}

\begin{solution}
设
\[
x=3+4t.
\]
代入第二个条件：
\[
3+4t\equiv 1\pmod 5
\Rightarrow 4t\equiv 3\pmod 5
\Rightarrow t\equiv 2\pmod 5.
\]
设
\[
t=2+5s,
\]
则
\[
x=11+20s.
\]
代入第三个条件：
\[
11+20s\equiv 6\pmod 7
\Rightarrow 6s\equiv 2\pmod 7.
\]
所以
\[
s\equiv 5\pmod 7.
\]
取最小 $s=5$，得
\[
\boxed{x=111}.
\]
\end{solution}

\begin{exercise}[B7]
写出斜边小于 $40$ 的全部原始勾股数。
\end{exercise}

\begin{solution}
枚举可得
\[
\boxed{(3,4,5),(5,12,13),(8,15,17),(7,24,25),(20,21,29),(12,35,37)}.
\]
\end{solution}

\begin{exercise}[B8]
求周长为 $60$ 的全部直角三角形整数边长。
\end{exercise}

\begin{solution}
由原始勾股数 $(3,4,5)$ 放大 $5$ 倍，得
\[
(15,20,25).
\]
由原始勾股数 $(5,12,13)$ 放大 $2$ 倍，得
\[
(10,24,26).
\]
故全部解为
\[
\boxed{(15,20,25),(10,24,26)}.
\]
\end{solution}

\begin{exercise}[B9]
求正整数解
\[
x^2-y^2=96.
\]
\end{exercise}

\begin{solution}
化为
\[
(x-y)(x+y)=96.
\]
枚举偶因数对：
\[
2\times 48,\quad 4\times 24,\quad 6\times 16,\quad 8\times 12.
\]
对应得到
\[
\boxed{(25,23),(14,10),(11,5),(10,2)}.
\]
\end{solution}

\begin{exercise}[B10]
判断方程
\[
x^2-y^2=2026
\]
是否有整数解，并说明理由。
\end{exercise}

\begin{solution}
因为
\[
2026\equiv 2\pmod 4,
\]
而平方差不可能模 $4$ 余 $2$，所以无整数解。
\end{solution}

\begin{exercise}[B11]
求正整数解
\[
xy=x+y+35.
\]
\end{exercise}

\begin{solution}
化为
\[
(x-1)(y-1)=36.
\]
枚举正因数对，可得
\[
\boxed{(2,37),(3,19),(4,13),(5,10),(7,7),(10,5),(13,4),(19,3),(37,2)}.
\]
\end{solution}

\begin{exercise}[B12]
求正整数解
\[
\frac1x+\frac1y=\frac18.
\]
\end{exercise}

\begin{solution}
化为
\[
(x-8)(y-8)=64.
\]
所以全部正整数解为
\[
\boxed{(9,72),(10,40),(12,24),(16,16),(24,12),(40,10),(72,9)}.
\]
\end{solution}

\begin{exercise}[B13]
求正整数解
\[
9x+14y=121.
\]
\end{exercise}

\begin{solution}
由
\[
9x\equiv 121\equiv 9\pmod{14}
\]
得
\[
x\equiv 1\pmod{14}.
\]
设
\[
x=1+14t,
\]
代回得
\[
y=8-9t.
\]
要使 $y>0$，只能取 $t=0$，所以唯一正整数解为
\[
\boxed{(x,y)=(1,8)}.
\]
\end{solution}

\begin{exercise}[B14]
证明：不存在整数 $x,y$，使得
\[
x^2+y^2=3.
\]
\end{exercise}

\begin{solution}
任意整数平方模 $4$ 只能余 $0$ 或 $1$，因此
\[
x^2+y^2
\]
模 $4$ 只能余 $0,1,2$，不可能余 $3$。故无整数解。
\end{solution}

\section{教学说明}

\begin{notebox}{教学建议}
\begin{itemize}
  \item 同余部分要让学生自己发现“循环节”和“分类讨论”的使用场景；
  \item 一次同余方程建议反复强调“先看最大公因数”；
  \item 韩信点兵先用逐步代入法，再介绍中国剩余定理；
  \item 勾股数部分不必一开始追求完整证明，但应逐步熟悉参数公式的使用；
  \item 不定方程部分重点培养“看结构、凑因式、做变形”的意识。
\end{itemize}
\end{notebox}

\begin{notebox}{课后加练建议}
可继续补充三类题：
\begin{enumerate}
  \item 多条件余数问题与韩信点兵变式；
  \item 周长、面积、斜边限制下的勾股数问题；
  \item 结合模运算证明无解的拓展题。
\end{enumerate}
\end{notebox}

\begin{landscape}
\newpage
\section{附加页：常用逆元速查表}

\begin{goalbox}{查表说明}
如果整数 $a$ 与模数 $m$ 互素，那么表中的“逆元”就是满足
\[
ax\equiv 1\pmod m
\]
的那个数。解一次同余方程时，如果模数比较小，可以直接查表使用。
\end{goalbox}

\begin{notebox}{查表提示}
\begin{itemize}
  \item 只有与模数互素的数才有逆元；
  \item 一个数的逆元不一定和它不同，例如模 $7$ 下的 $6$ 的逆元还是 $6$；
  \item 如果 $a$ 的逆元是 $b$，那么 $b$ 的逆元也是 $a$。
\end{itemize}
\end{notebox}

\subsection*{横向速查表}

\[
\renewcommand{\arraystretch}{1.35}
\begin{array}{c|cccccccccccc}
\text{模数} & 1 & 2 & 3 & 4 & 5 & 6 & 7 & 8 & 9 & 10 & 11 & 12 \\
\hline
7  & 1 & 4 & 5 & 2 & 3 & 6 & - & - & - & - & - & - \\
9  & 1 & 5 & - & 7 & 2 & - & 4 & 8 & - & - & - & - \\
11 & 1 & 6 & 4 & 3 & 9 & 2 & 8 & 7 & 5 & 10 & - & - \\
12 & 1 & - & - & - & 5 & - & 7 & - & - & - & 11 & - \\
13 & 1 & 7 & 9 & 10 & 8 & 11 & 2 & 5 & 3 & 4 & 6 & 12
\end{array}
\]

\begin{notebox}{读表示例}
例如在“模 $11$”这一行、数字 $7$ 这一列，查到的结果是 $8$，表示
\[
7^{-1}\equiv 8\pmod{11}.
\]
横线 “$-$” 表示该数与模数不互素，因此不存在逆元。
\end{notebox}

\begin{notebox}{记忆提示}
这些表不要求死记硬背，但下面几组建议优先熟悉：
\[
2^{-1}\equiv 4\pmod 7,\qquad
3^{-1}\equiv 5\pmod 7,\qquad
2^{-1}\equiv 5\pmod 9,
\]
\[
3^{-1}\equiv 4\pmod{11},\qquad
7^{-1}\equiv 8\pmod{11},\qquad
4^{-1}\equiv 10\pmod{13}.
\]
这些在竞赛题中出现得很频繁。
\end{notebox}
\end{landscape}

\newpage
\section{附加页：常用平方数模表}

\begin{goalbox}{平方数模表的用途}
很多题目需要判断一个式子是否可能等于某个平方数，或者判断方程是否无解。
这时只要先看“平方数除以某个数后可能余什么”，常常就能迅速排除不可能的情况。
\end{goalbox}

\begin{notebox}{使用说明}
例如模 $8$ 的平方数只可能余 $0,1,4$，所以如果某个数模 $8$ 余 $3$、$5$、$6$ 或 $7$，它就不可能是平方数。
\end{notebox}

\subsection*{常用平方数模表}

\[
\renewcommand{\arraystretch}{1.4}
\begin{array}{c|c}
\text{模数} & \text{平方数可能的余数} \\
\hline
4 & 0,\,1 \\
8 & 0,\,1,\,4 \\
3 & 0,\,1 \\
5 & 0,\,1,\,4 \\
7 & 0,\,1,\,2,\,4 \\
9 & 0,\,1,\,4,\,7
\end{array}
\]

\subsection*{从哪里来：以模 $7$ 为例}

只要把模 $7$ 的所有可能余数平方一遍即可：
\[
0^2\equiv 0,\quad 1^2\equiv 1,\quad 2^2\equiv 4,\quad 3^2\equiv 2\pmod 7.
\]
后面的
\[
4^2,\ 5^2,\ 6^2
\]
分别与前面重复，所以平方数模 $7$ 只可能余
\[
0,1,2,4.
\]

\begin{notebox}{记忆提示}
下面几条最常用，建议优先记熟：
\begin{itemize}
  \item 平方数模 $4$ 只能余 $0$ 或 $1$；
  \item 平方数模 $8$ 只能余 $0,1,4$；
  \item 平方数模 $3$ 只能余 $0$ 或 $1$；
  \item 平方数模 $5$ 只能余 $0,1,4$。
\end{itemize}
这几条在证明“不可能是平方数”时特别高频。
\end{notebox}

\end{document}
